\documentclass[10pt,a4paper]{amsart}
\usepackage[margin=19mm]{geometry}
\usepackage{amsmath,amssymb,mathtools}
\usepackage[scr=rsfs,cal=euler]{mathalfa}
\usepackage{xfrac}
\usepackage[unicode,hidelinks,bookmarks=false]{hyperref}
\newcommand{\ie}{i.e.\ }
\newcommand{\cf}{cf.\ }
\newcommand{\resp}{respectively\ }
\newcommand{\Subsection}{\subsection}
\providecommand{\nobreakdash}{\nobreak\hbox{-}}
\setlength{\emergencystretch}{2em}
\multlinegap0pt
\usepackage{amssymb}
\usepackage{bm}
\usepackage{tikz}
\usepackage[all]{xy}
\newtheorem{prop}{Proposition}
\newtheorem{theorem}{Theorem}
\DeclareMathOperator{\SP}{SP}
\DeclareMathOperator{\card}{card}
\newcommand{\sci}{\wedge}
\DeclareMathOperator{\supp}{supp}
\DeclareMathOperator{\vol}{vol}
\title{Overlapping substitutions and tilings}
\author{Shigeki Akiyama, Yasushi Nagai, Shu-Qin Zhang}
\date{Source: arXiv:2407.18666v3 (CC BY 4.0)}
\begin{document}
\maketitle

\noindent\emph{Source and licence.} Overlapping substitutions and tilings. Version arXiv:2407.18666v3 (CC BY 4.0), distributed by arXiv under the Creative Commons Attribution 4.0 International licence, http://creativecommons.org/licenses/by/4.0/, as recorded in the arXiv licence metadata for this exact version and shown on the abstract page https://arxiv.org/abs/2407.18666v3. Full licence notice: \texttt{licenses/arxiv-2407.18666-CC-BY-4.0.txt}.\par\smallskip

\noindent\textit{Editorial scope.} Counted unit: the article's theorem thm\_uniform\_patch\_freq (source0.tex lines 920--932). The article's complete original proof of it is delivered with it (source0.tex lines 934--995, one \texttt{proof} environment of 62 lines). No mathematics is written, repaired or re-derived: the statement and the proof are the article's own lines, in the article's own notation, and no retained line is substituted or re-flowed.  Retained uncounted as the source-local context the proof needs: lines 758--785; lines 905--915.  Nothing was omitted from any retained span, so the package carries no omission record.  No retained line contains a citation, so the wrapper carries no bibliography and no bibliography entry.  No retained line contains a figure environment, an image-inclusion command or drawing code, so no image is delivered and the package declares no asset dependency.  Editorial formatting only, in addition: an AMS \texttt{amsart} preamble that restates the article's own declarations and macros used by the retained lines, namely the environments \texttt{prop}, \texttt{theorem} on separate counters, together with the standard packages those lines need.  The retained environments print in this document as Proposition 1, Theorem 1 in order of appearance, whereas the article prints the same items with the numbering of its own full document.  The complete native source of the article as distributed in the arXiv e-print for this exact version is bundled unchanged as \texttt{sources/162926/source0.tex}.
\begin{prop}\label{prop_frequency_weighted-substi}
 Let $v$ be such that
\begin{enumerate}
\item $\SP(v)$ is a tiling in $\mathbb{R}^d$,
 \item $v(T)=1$ for any $T\in\SP(v)$,
\item $\xi^k(v)=v$ for some $k>0$, and
\item the substitution matrix for $v$ is primitive.
\end{enumerate}
Take a Perron--Frobenius left eigenvector $l=(l_1,l_2,\ldots ,l_{\kappa})$ and
right eigenvector
\begin{align*}
 r=
\begin{pmatrix}
 r_1\\r_2\\\vdots\\r_{\kappa}
\end{pmatrix}
\end{align*}
for the substitution matrix for $\xi$, 
and assume $l_i$ coincides with the Lebesgue measure $\vol(\supp T_i)$
for $\supp T_i$ for each $i$, and
$l$ and $r$ are normalized in the sense that $\sum_{i=1}^{\kappa}l_ir_i=\sum_{i=1}^{\kappa}r_i=1$.

Then $r$ describes the frequency for $v$, which is regarded as a tiling: for any van Hove sequence 
$(A_n)_{n=1,2,\ldots}$ and $x_1,x_2,\ldots\in\mathbb{R}^d$, we have a convergence
\begin{align*}
 \lim_{n\rightarrow \infty}\frac{1}{\vol(A_n)}\tau_{T_i}(v\sci (A_n+x_n))=r_i,
\end{align*}
which is uniform for $(x_n)$.
\end{prop}

\begin{theorem}\label{thm_uniform_tile_freq}
 Let $\rho$ be an expanding overlapping subsitutiion and $r=(r_1,r_2,\ldots, r_{\kappa})$ be
a right Perron--Frobenius eigenvector which is normalized in the sense of 
Proposition \ref{prop_frequency_weighted-substi}. Then $r$ gives the frequency for any fixed points for $\rho$.
In particular, for any van Hove sequence $(A_n)$ and elements $x_n\in\mathbb{R}^d$, we have
a convergence
\begin{align*}
      \lim_{n\rightarrow\infty}\frac{\card\{x\in A_n+x_n\mid T_i+x\in\mathcal{T}\}}{\vol{A_n}}=r_i,
\end{align*}
which is uniform for $(x_n)$.
\end{theorem}

\begin{theorem}\label{thm_uniform_patch_freq}
    Assume a fixed point $\mathcal{{T}}$ for $\rho$ is repetitive. Then the patch frequencies converge uniformly. That is, 
    for any van Hove sequence $(A_n)$,
    elements $x_n\in\mathbb{R}^d$
    and a patch $\mathcal{{P}}$, the
    limit
    \begin{align*}
    \lim_{n\rightarrow\infty}\frac{\card\{x\in A_n+x_n\mid \mathcal{P}+x\subset\mathcal{T}\}}{\vol A_n}
    \end{align*}
    converges and is independent of $(A_n)$ and $x_n$.
    Moreover, if we fix $(A_n)$ and $\mathcal{P}$, the limit is uniform for $(x_n)$.
    In particular, the corresponding dynamical system for the repetitive fixed point is uniquely ergodic.
\end{theorem}

\begin{proof}
    For a $P\in\mathcal{B}$
    and an $x\in\mathbb{R}^d$, we define the representative point for $P+x$ via
    \begin{align*}
        c(P+x)=x.
    \end{align*}
    We may assume that $0\in\mathbb{R}^d$ is in the interior of each prototile $P\in\mathcal{{B}}$.

    For each $L>0$, we
    define an additional label for tiles in
    $\mathcal{T}$ by 
    looking at the patterns around the tiles inside the ball of raduis $L$.
    For  a $T\in\mathcal{T}$, 
    we will look at
\begin{align*}\mathcal{T}\sqcap B(c(T),L),
    \end{align*}
    but we want to identify two such patches that are translationally equivalent.
    The
    corresponding ``coloured'' tile $\tilde{T}$ is the tile $T$ with additional label or colour
    \begin{align*}
        (\mathcal{T}-c(T))\sqcap B_L.
    \end{align*}

    The new alphabet for the coloured tiling is
    \begin{align*}
        \tilde{\mathcal{B}}=\{(T-c(T),(\mathcal{T}-c(T))\sqcap B_L)\mid T\in\mathcal{T}\}
        =\{\tilde{T}-c(T)\mid T\in\mathcal{T}\}.
    \end{align*} 
    We define a substitution $\tilde{\rho}$ via
    \begin{align*}
        \tilde{\rho}(\tilde{T}-c(T))=
        \{\tilde{S}-\phi(c(T))\mid S\in\rho(T)\}.
    \end{align*}
    This is well-defined.
    This means that the image for $\tilde{T}$ is
    \begin{align*}
        \tilde{\rho}(\tilde{T})=\{\tilde{S},\mid S\in\rho(T)\}=\{(S,(\mathcal{T}-c(S))\sqcap B_L)\mid S\in\rho(T)\}.
    \end{align*}
    We set the colourd tiling via
    \begin{align*}
        \tilde{\mathcal{T}}=\{\tilde{T}\mid T\in\mathcal{T}\}.
    \end{align*}
    It is easy to show that $\tilde{\rho}(\tilde{\mathcal{T}})=\tilde{\mathcal{T}}$.

    For a patch $\mathcal{P}$, we take an arbitrary $S\in\mathcal{P}$ and fix it. We set
    \begin{align*}
        \widetilde{\mathcal{B}}_{\mathcal{P}}
        =\{(P,\mathcal{S})\in\tilde{\mathcal{B}}\mid \mathcal{P}-c(S)\subset \mathcal{S}\}.
    \end{align*}
        We have a bijection
        \begin{align*}
            &x\in\{x\in\mathbb{R}^d\mid \mathcal{P}+x\subset\mathcal{T}\}\\
            &\mapsto ((P,(\mathcal{T}-x-c(S))\sqcap B_L), x+c(S))\in\{((P,\mathcal{S}),y)\in\widetilde{\mathcal{B}}_{\mathcal{P}}\times \mathbb{R}^d\mid (P,\mathcal{S})+y\in\tilde{\mathcal{T}}\}.
        \end{align*}
We conclude
\begin{align*}
    \card\{x\in A_n\mid \mathcal{P}+x\subset\mathcal{T}\}=\sum_{\tilde{P}\in\widetilde{\mathcal{B}}_{\mathcal{P}}}\card\{x\in A_n\mid \tilde{P}+x\in\tilde{\mathcal{T}}\}+o(\vol (A_n)).
\end{align*}
The repetitivity of $\mathcal{T}$ implies the primitivity of $\tilde{\rho}$, which implies the uniform convergence for the tile frequencies for $\tilde{\mathcal{T}}$ by
Theorem \ref{thm_uniform_tile_freq}.
The theorem follows.
\end{proof}

\end{document}
